\item \textbf{Indica qué diagramas representan la información requerida por las siguientes solicitudes de información:}

\textbf{¿Qué taxis maneja el chofer Raúl López en el sitio Santa Fe?}

Solo los diagramas a) y c) dejan responder esto. Ambos capturan en una relación el hecho de que un chofer maneja un taxi específico en un sitio específico (a) mediante una relación ternaria y c) mediante la combinación de Tener con Manejar), por lo que se puede filtrar directamente por chofer y sitio para obtener el taxi. En el diagrama b), Manejar sólo relaciona Chofer con Taxi, sin ninguna referencia al sitio, así que no podemos saber en qué sitio maneja Raúl López un taxi específico (básicamente la situación que sale en errores comunes en las notas).

\textbf{¿A qué sitios está afiliado el chofer Carlos Reyes?}

De nuevo, solo a) y c) lo permiten, ya que en ambos el chofer queda relacionado al sitio con una misma relación. En b) el chofer solo se relaciona con los taxis que maneja, y el sitio se obtiene por separado a través de Tener(Sitio, Taxi) y como no hay garantía de que el chofer maneje ese taxi específicamente en ese sitio (el taxi podría estar en ese sitio en otro momento o con otro chofer), no se puede inferir que un chofer pertenece a un sitio.

\textbf{¿Qué taxis están asociados al sitio Universidad?}

Esta pregunta la responden los tres diagramas, en todos existe una relación (directa en b y c, por proyección en a) entre Sitio y Taxi que no depende del chofer.

\item \textbf{¿Qué modificación harías en el diagrama de la figura a), sin perder información, para que se puedan conocer qué taxis que maneja cada chofer?}

Agregaría una relación binaria adicional \textit{Maneja} directamente entre Chofer y Taxi, independiente de la relación \textit{Tener}. De esta forma, conservamos toda la información original (qué chofer maneja qué taxi en qué sitio usando \textit{Tener}) y se puede consultar directamente qué taxis maneja un chofer sin importar el sitio.

\item \textbf{¿Qué diferencia existe entre los diagramas de las figuras a) y c)?}

Ambos diagramas terminan representando la misma información (qué chofer maneja qué taxi en qué sitio), pero lo hacen de forma distinta. En a) se usa una sola relación de grado tres (ternaria) que conecta Sitio, Chofer y Taxi. En c) se usa una agregación, primero se define la relación binaria \textit{Tener} entre Sitio y Taxi y después esa relación completa se trata como si fuera una entidad para relacionarla usando \textit{Manejar} con Chofer. en otras palabras, c) descompone en dos pasos binarios lo que a) resuelve en un paso, pero ambos preservan la combinación de los tres elementos a la vez, a diferencia de b).

\item \textbf{Cómo modificarías el diagrama de la figura a) para representar las siguientes restricciones:}

\textbf{Un chofer no puede manejar más de un taxi en el mismo sitio.}

Se agregaría una restricción de cardinalidad sobre la relación ternaria \textit{Tener}, colocando un 1 del lado de Taxi. Esto se interpreta como que, para cada combinación de (Sitio, Chofer), existe a lo más un Taxi asociado.

\textbf{Un taxi no puede ser manejado por más de un chofer en el mismo sitio.}

De manera análoga, se coloca un 1 del lado de Chofer, para indicar que por cada combinación de (Sitio, Taxi) existe a lo más un Chofer asociado.
